\documentclass[letterpaper,10pt]{article}

\usepackage[left=0.6in,top=0.6in,right=0.6in,bottom=0.6in]{geometry}
\usepackage{hyperref}
\usepackage{enumitem}
\usepackage{titlesec}
\usepackage{xcolor}
\usepackage{tabularx}

% Page setup
\pagestyle{empty}
\urlstyle{same}
\raggedbottom
\raggedright
\setlength{\tabcolsep}{0in}

% Colors
\definecolor{primaryDark}{RGB}{15, 23, 42}
\definecolor{accentBlue}{RGB}{37, 99, 235}
\definecolor{textGray}{RGB}{71, 85, 105}
\definecolor{ruleGray}{RGB}{203, 213, 225}

% Hyperref setup
\hypersetup{
    colorlinks=true,
    linkcolor=accentBlue,
    urlcolor=accentBlue,
    pdftitle={Varizky Naldiba Rimra - Curriculum Vitae},
    pdfauthor={Varizky Naldiba Rimra}
}

% Section formatting
\titleformat{\section}{
  \vspace{-4pt}\raggedright\large\bfseries\color{primaryDark}
}{}{0em}{}[\color{ruleGray}\titlerule \vspace{-3pt}]

% Custom entry macros
\newcommand{\resumeItem}[1]{
  \item\small{
    {#1 \vspace{-1.5pt}}
  }
}

\newcommand{\resumeSubheading}[4]{
  \vspace{-1pt}\item
    \begin{tabularx}{\linewidth}[t]{@{}X r@{}}
      \textbf{#1} & \textbf{#2} \\
      \textit{\small#3} & \textit{\small #4} \\
    \end{tabularx}\vspace{-4pt}
}

\newcommand{\resumeProjectHeading}[2]{
  \vspace{-1pt}\item
    \begin{tabularx}{\linewidth}[t]{@{}X r@{}}
      #1 & \textit{\small #2} \\
    \end{tabularx}\vspace{-4pt}
}

\newcommand{\resumeSubHeadingListStart}{\begin{itemize}[leftmargin=0.05in, label={}]}
\newcommand{\resumeSubHeadingListEnd}{\end{itemize}}
\newcommand{\resumeItemListStart}{\begin{itemize}[leftmargin=0.15in, itemsep=1.5pt]}
\newcommand{\resumeItemListEnd}{\end{itemize}\vspace{-3pt}}

\begin{document}

%----------HEADER----------
\begin{center}
    {\Huge \bfseries \color{primaryDark} Varizky Naldiba Rimra} \\ \vspace{3pt}
    \small \textbf{Applied Machine Learning \& IoT Systems Engineer} \\ \vspace{4pt}
    \small Malang, East Java, Indonesia \ $|$ \ 
    \href{mailto:varizkynr@gmail.com}{varizkynr@gmail.com} \ $|$ \ 
    \href{https://linkedin.com/in/varizkynaldiba/}{linkedin.com/in/varizkynaldiba} \ $|$ \ 
    \href{https://github.com/VarizkyNaldiba}{github.com/VarizkyNaldiba}
\end{center}
\vspace{-6pt}

%-----------PROFESSIONAL SUMMARY-----------
\section{Professional Summary}
\begin{flushleft}
\small
Informatics Engineer from Politeknik Negeri Malang specializing in Applied Machine Learning, IoT Sensor Telemetry, Computer Vision pipelines, and Cloud-Connected Data Systems. Proven track record in designing resilient systems from edge hardware and sensor calibration curves (ESP32/Arduino) to deep-learning crowd density estimation models, regional educational accessibility analyses, and full-stack decision-support platforms.
\end{flushleft}
\vspace{-6pt}

%-----------EDUCATION-----------
\section{Education}
  \resumeSubHeadingListStart
    \resumeSubheading
      {Politeknik Negeri Malang (State Polytechnic of Malang)}{Malang, East Java, Indonesia}
      {Bachelor of Applied Science (D4) in Informatics Engineering}{2022 -- Present}
      \resumeItemListStart
        \resumeItem{Focus Areas: Applied Artificial Intelligence, Embedded Systems \& IoT Telemetry, Geospatial Analytics, Systems Architecture.}
        \resumeItem{Project-Based Learning: Multi-tier compliance systems, academic achievement matching platforms, and telemetry calibration.}
      \resumeItemListEnd
  \resumeSubHeadingListEnd

%-----------TECHNICAL SKILLS-----------
\section{Technical Skills}
 \begin{itemize}[leftmargin=0.05in, label={}]
    \small{\item{
     \textbf{Core Languages}{: Python, C/C++, PHP, JavaScript, SQL, HTML5, Modern CSS} \\ \vspace{1pt}
     \textbf{Machine Learning \& AI}{: TensorFlow, Keras, OpenCV, Scikit-Learn, LSTM Networks, CNNs (MCNN/CSRNet), Support Vector Machines (SVM), GLCM \& HOG Feature Extraction, Google Gemini API, Prompt Engineering} \\ \vspace{1pt}
     \textbf{Data Science \& Geospatial}{: Pandas, NumPy, Matplotlib, GeoJSON, Folium GIS, Spatial Disparity Indexing, Levenshtein Fuzzy Matching, Statistical Regression, ETL Pipelines} \\ \vspace{1pt}
     \textbf{Web \& Systems Architecture}{: React, Next.js, Laravel, Streamlit, Firebase Realtime Database, MySQL, RESTful APIs, Decision Support Systems (DSS), Bootstrap} \\ \vspace{1pt}
     \textbf{IoT \& Embedded Systems}{: ESP32, Arduino, Telemetry Probes (pH, Turbidity, DS18B20 Temp, DO), Non-linear Calibration Curves, Local Flash Ring Buffer Failover} \\ \vspace{1pt}
     \textbf{Engineering Tools \& Platforms}{: Git, GitHub, Figma, HuggingFace Spaces, Google Colab, VS Code, Linux}
    }}
 \end{itemize}
 \vspace{-8pt}

%-----------ENGINEERED PROJECTS-----------
\section{Engineered Projects}
  \resumeSubHeadingListStart

    \resumeProjectHeading
      {\textbf{Aquavion — IoT Aquaculture Telemetry Platform} $|$ \emph{Next.js, Firebase, ESP32, Arduino, C++}}{\href{https://aqua-vion.vercel.app/}{aqua-vion.vercel.app}}
      \resumeItemListStart
        \resumeItem{Architected an industrial IoT telemetry monitoring platform for freshwater catfish aquaculture, streaming real-time sensor parameters (pH, DS18B20 temperature, nephelometric turbidity) to prevent fish mortality.}
        \resumeItem{Engineered microcontroller firmware with non-linear polynomial calibration offset curves and automated temperature compensation for continuous submerged probe stability.}
        \resumeItem{Implemented a non-volatile circular flash ring buffer on the ESP32, guaranteeing zero telemetry data loss during agricultural network downtime by caching readings and auto-synchronizing upon reconnection.}
        \resumeItem{Built a responsive Next.js analytics web dashboard synchronized via Firebase Realtime Database with threshold-based visual alerts and telemetry trend charting.}
      \resumeItemListEnd

    \resumeProjectHeading
      {\textbf{AI Restaurant Advisory System} $|$ \emph{Python, Google Gemini API, Streamlit, Prompt Engineering}}{\href{https://varizkynaldiba-project-of-llm-based-tools-and-gem-fnb-rm-vjbmhr.streamlit.app/}{Live Web App}}
      \resumeItemListStart
        \resumeItem{Developed an LLM-backed decision-support system for food service enterprises (RM Padang) synthesizing financial margins, seasonal ingredient cost swings, and menu item demand.}
        \resumeItem{Designed structured prompt engineering schemas with domain context injection, enforcing deterministic calculations for dish cost margins and shelf-life metrics while eliminating hallucinations.}
        \resumeItem{Built an interactive Streamlit analytics interface incorporating an automated BCG menu engineering matrix to classify dishes by profitability versus popularity.}
      \resumeItemListEnd

    \resumeProjectHeading
      {\textbf{Fruit Classification Computer Vision Pipeline} $|$ \emph{Python, OpenCV, Scikit-Learn, HuggingFace}}{\href{https://huggingface.co/spaces/Zelectra/PBL.MachineLearningAndPCVK}{HuggingFace Space}}
      \resumeItemListStart
        \resumeItem{Engineered an invariant multi-class image recognition pipeline classifying 16 fruit varieties using hybrid feature extraction coupled with hyperparameter-tuned Support Vector Machines (SVM).}
        \resumeItem{Constructed a multi-scale feature pipeline combining Gray-Level Co-occurrence Matrix (GLCM texture), 2D Discrete Wavelet Transform, and Histogram of Oriented Gradients (HOG edge gradients).}
        \resumeItem{Applied Contrast Limited Adaptive Histogram Equalization (CLAHE) and foreground segmentation masks to achieve illumination-invariant feature vectors; deployed real-time inference space on HuggingFace.}
      \resumeItemListEnd

    \resumeProjectHeading
      {\textbf{SLB Spatial Disparity \& Accessibility Analysis} $|$ \emph{Python, Pandas, GeoJSON, Folium}}{\href{https://huggingface.co/spaces/Tim-Maju-Lancar-Datathon/Dashboard-Tim-Maju-Lancar-Datathon-Jatim-2025}{Interactive Dashboard}}
      \resumeItemListStart
        \resumeItem{Conducted geodemographic spatial analysis across 38 East Java regencies to evaluate accessibility and infrastructure equity for Special Education Schools (SLB) for Datathon Jatim 2025 (Team Tim Maju Lancar).}
        \resumeItem{Engineered an automated Python ETL pipeline leveraging Levenshtein fuzzy string matching to harmonize disparate municipal school registries and bind attributes directly to GeoJSON polygon boundaries.}
        \resumeItem{Formulated a Spatial Accessibility Index (SAI) correlating school capacity with regional youth disability statistics, deploying interactive choropleth dashboards on HuggingFace Spaces.}
      \resumeItemListEnd

    \resumeProjectHeading
      {\textbf{Crowd Density Estimation Model} $|$ \emph{Python, TensorFlow, Keras, CNNs, OpenCV}}{\href{https://github.com/VarizkyNaldiba/PenyelamHandal_Hology}{GitHub Repository}}
      \resumeItemListStart
        \resumeItem{Designed a deep convolutional density map regression system for the Hology Data Mining competition (Team Penyelam Handal) to estimate crowd counts in unconstrained surveillance imagery.}
        \resumeItem{Implemented a Multi-Column Convolutional Neural Network (MCNN) architecture with parallel small (3x3), medium (5x5), and large (7x7) filter branches to extract scale-invariant density features under heavy occlusions.}
        \resumeItem{Trained and augmented across 1,900 benchmark training images; validated model convergence against 500 test images using Mean Absolute Error (MAE) and Mean Squared Error (MSE) metrics.}
      \resumeItemListEnd

    \resumeProjectHeading
      {\textbf{AlgaSense — Marine Seaweed Aquaculture IoT \& ML} $|$ \emph{Python, TensorFlow, Keras, LSTM}}{\href{https://github.com/VarizkyNaldiba/Algasense}{GitHub Repository}}
      \resumeItemListStart
        \resumeItem{Engineered subsea telemetry monitoring and deep learning growth forecasting for marine seaweed cultivators in Sumenep, Madura, using LSTM neural networks to forecast optimal harvest cycles.}
        \resumeItem{Trained Long Short-Term Memory recurrent networks on continuous oceanic telemetry (salinity, pH, temperature, dissolved oxygen) with moving-window noise filtration to forecast growth trajectories.}
      \resumeItemListEnd

    \resumeProjectHeading
      {\textbf{SIMPRES — Student Achievement Management System} $|$ \emph{Laravel, PHP, MySQL, Bootstrap, DSS}}{\href{https://github.com/VarizkyNaldiba/Prestasi-Mahasiswa-S4}{GitHub Repository}}
      \resumeItemListStart
        \resumeItem{Developed an enterprise academic portal for the Information Technology Department at Politeknik Negeri Malang to streamline student competition verification and mentorship tracking.}
        \resumeItem{Integrated a Decision Support System (DSS) matching student skill profiles to national/international contests, alongside normalized MySQL schemas with transactional ACID guarantees.}
      \resumeItemListEnd

  \resumeSubHeadingListEnd

%-----------HONORS & ACHIEVEMENTS-----------
\section{Honors \& Verified Competitions}
  \resumeSubHeadingListStart
    \resumeSubheading
      {TFS 2026 LKTI Nasional — Juara Harapan 2}{Malang, Indonesia}
      {CatfishCare IoT Aquaculture Telemetry — SK No. 27881/PL2.3/KM/2026}{2026}
      \resumeItemListStart
        \resumeItem{Awarded 2nd Runner Up (Harapan 2) at national scientific writing contest for resilient embedded IoT aquaculture system.}
      \resumeItemListEnd
    \resumeSubheading
      {PKM POLINEMA 2026 — Lolos Seleksi Tahap 2 Insentif}{Malang, Indonesia}
      {Smart Aquaculture Integration System (SAIS) — SK No. 40/WADIR.III/KM/2026}{2026}
      \resumeItemListStart
        \resumeItem{Selected for Phase 2 Incentive Scheme funding for agricultural IoT telemetry and failover architecture.}
      \resumeItemListEnd
    \resumeSubheading
      {Intelecta Cup 2025 — Peringkat 18 dari 82 Tim Nasional}{Indonesia}
      {Team P-1 ML — Agroclimatology Time-Series Weather Forecasting}{2025}
      \resumeItemListStart
        \resumeItem{Achieved Kaggle competition score of 94.310, securing 18th position nationwide among 82 competing university teams.}
      \resumeItemListEnd
    \resumeSubheading
      {TFS 2025 LKTI Nasional — Finalis Nasional}{Malang, Indonesia}
      {AlgaSense Marine Seaweed Aquaculture IoT \& LSTM — SK No. 16002/PL2.3/KM/2025}{2025}
      \resumeItemListStart
        \resumeItem{Selected as national finalist for ocean telemetry and LSTM neural network predictive harvest modeling.}
      \resumeItemListEnd
    \resumeSubheading
      {Hology 8.0 Data Mining Competition — Official Participant}{Malang, Indonesia}
      {Team Penyelam Handal — Multi-Column CNN Crowd Density Estimation}{2025}
      \resumeItemListStart
        \resumeItem{Engineered multi-column convolutional density regression models for unconstrained surveillance imagery at national data mining contest.}
      \resumeItemListEnd
    \resumeSubheading
      {Jatim Datathon 2025 — Official Participant (Peserta Resmi)}{Surabaya, Indonesia}
      {Team Tim Maju Lancar — Spatial Disparity Analysis (Diskominfo Jawa Timur)}{2025}
      \resumeItemListStart
        \resumeItem{Conducted geodemographic accessibility analysis for Special Education Schools across 38 East Java regencies.}
      \resumeItemListEnd
  \resumeSubHeadingListEnd

%-----------DIGITAL MEDIA & PRODUCTION-----------
\section{Digital Media \& Production}
  \resumeSubHeadingListStart
    \resumeProjectHeading
      {\textbf{Content Producer \& Video Editor} $|$ \emph{Motion Graphics, Video Pacing, Community Analytics}}{2022 -- Present}
      \resumeItemListStart
        \resumeItem{Produce short-form and long-form video game analysis, competitive mechanics guides, and motion graphics reels across TikTok (\href{https://www.tiktok.com/@oryzasatirva}{@oryzasatirva}) and Bilibili (\href{https://www.bilibili.tv/id/space/1416802109}{Space 1416802109}).}
        \resumeItem{Analyze audience retention curves and distribution algorithms, optimizing editorial pacing and thumbnail engagement.}
      \resumeItemListEnd
  \resumeSubHeadingListEnd

\end{document}
